
\subsubsection{6.1 Findings}

\textbf{Finding 1: Registry assembly is feasible and reveals ecosystem documentation patterns.} 96.5\% of Open LLM Leaderboard models document zero curation practices. The 3.5\% that do are concentrated in well-established model families. This sparsity characterizes the documentation landscape within which AI governance frameworks must operate.

\textbf{Finding 2: Naive regression conflates documentation with model characteristics.} The negative beta\_docs = -3.45 is statistically robust but substantively misleading as a causal estimate. It reflects a selection artifact: documented models are from an earlier, smaller-parameter era. Without size-matched controls, the documentation-quality relationship cannot be credibly estimated from these data.

\textbf{Finding 3: Perplexity filtering is undocumented under its canonical name.} Complete absence of the \texttt{perplexity\textbackslash \_filter\textbackslash \_documented} signal across 3,749 model cards indicates that automated documentation checks using canonical vocabulary will systematically miss this practice. Labs use alternative terms such as ''quality filtering,'' ''CCNet filtering,'' or ''fastText quality filter.''

\subsubsection{6.2 Limitations}

\textbf{L1: Documentation-to-implementation validity is untested.} The assumption that documentation presence reflects actual curation implementation has not been verified. Bridge validation using models with accessible pretraining corpora (e.g., Pythia/The Pile, OLMo/Dolma) is required.

\textbf{L2: Extreme feature sparsity.} With 96.5\% of models at doc\_score = 0, the effective sample of documented models is small (n = 158). This is sufficient for the existence gate but may limit statistical power for subgroup analyses.

\textbf{L3: V2 benchmark mismatch with original hypothesis.} The hypothesis was originally formulated for V1 knowledge-recall benchmarks (MMLU, ARC). The operational data uses V2 benchmarks (IFEval, BBH, MATH Lvl 5, GPQA, MUSR, MMLU-PRO), which emphasize reasoning and mathematics. The original prediction about differential sensitivity across benchmark types requires reformulation for V2.

\textbf{L4: Observational design.} All beta\_docs estimates are subject to unmeasured confounding. Findings are correlational.

\textbf{L5: Targeted sampling introduces selection bias.} The families chosen (LLaMA, Mistral, Qwen, Falcon, Pythia, OLMo) are among the best-documented in the ecosystem. This non-random sampling is necessary for achieving feature variance but limits generalizability.

\textbf{L6: Missing card non-randomness.} The 744 models without retrievable cards are assigned doc\_score = 0 by default. These models may differ systematically from those with accessible cards, introducing measurement error that is not random with respect to benchmark performance.

\textbf{L7: Binary indicator operationalization.} The \texttt{perplexity\textbackslash \_filter\textbackslash \_documented} indicator uses canonical terminology and misses equivalent practices described under alternative names. This operationalization limitation applies to all four binary features and should be addressed in future work using vocabulary-aware NLP extraction.

\textbf{L8: Log(doc\_score + 1) specification not empirically verified.} An alternative log-transformed specification was considered to address the right-skewed distribution of doc\_score but was not computed from actual data. The directional consistency of such a specification with the linear model remains to be confirmed empirically.

\subsubsection{6.3 Broader Impact}

The registry provides a public dataset for empirically grounding AI governance frameworks. The finding that perplexity filtering documentation is absent under its canonical name has practical value for documentation standard bodies.

Uncritical use of this characterization could stigmatize model families without accounting for the documented selection biases. The negative beta\_docs should not be interpreted as evidence that documentation is harmful.
